\section{Introduction}
\label{sec:intro}

The question ``is this map partisan?'' has become central to redistricting
litigation.  Ensemble methods provide one framework for answering it: compare
the map to the distribution of plans consistent with geographic constraints,
and ask whether the map falls in an extreme tail.  But this framework is
typically applied to human-drawn maps that may have been deliberately designed
to achieve partisan ends.  How should it be applied to an algorithmic map
whose objective function has no partisan component?

This paper addresses that question for \ar{} \citep{dellaLibera2026b11}.
The \ar{} algorithm minimises edge cuts in the census-tract adjacency graph
under the prime factorization of the seat count.  No electoral data enters.
No partisan registration data enters.  No incumbency data enters.  The
algorithm is politically blind by construction.

Yet political blindness does not guarantee partisan neutrality.  An algorithm
that minimises edge cuts will tend to produce compact, geographically cohesive
districts.  Compact districts in states with urban-rural partisan sorting will
tend to contain urban Democratic majorities and rural Republican majorities,
potentially producing a Republican seat advantage even without any partisan
intent.  This is the Rodden effect \citep{rodden2019why}: the geographic
concentration of Democrats in cities is the root cause of the partisan gap
between votes and seats under any compact map.

The central question is therefore not whether \ar{} is partisan-blind
(it is) but whether its outputs are partisan-neutral in the sense of
producing outcomes within the range of what geography permits.  The ensemble
framework provides the answer.

\paragraph{Main findings.}
\begin{enumerate}
  \item \ar{} produces partisan outcomes within one seat of the \recom{}
    ensemble median in five of six studied states.  The median percentile
    across states is the 53rd percentile of Democratic seats.
  \item Georgia is the exception: the \ar{} plan produces 5D/9R at the
    38th percentile, reflecting strong geographic sorting, not algorithmic
    bias.
  \item The ``proportionality corridor'' (fraction of ensemble plans within
    one seat of proportional) is 20--40\% for competitive states.  The
    \ar{} plan falls inside this corridor for NC, WI, and MN (a state
    not studied in G.1 but analysed here for the corridor metric), but
    not for Georgia.
  \item Among three competing algorithms studied in B.0, only \ar{} and
    the flip-chain-optimal algorithm produce consistently median partisan
    outcomes; the short-burst algorithm can produce extreme outcomes if
    its objective is partisan-correlated.
\end{enumerate}

\paragraph{Organisation.}
Section~\ref{sec:bell-curve} describes the partisan distribution of ensemble
plans.  Section~\ref{sec:ar-position} quantifies the \ar{} plan's position.
Section~\ref{sec:corridor} defines and computes the proportionality corridor.
Section~\ref{sec:algorithm-choice} explains why algorithm choice determines
partisan tail.  Section~\ref{sec:conclusion} concludes.
